\documentclass[11pt,a4paper]{article}

\usepackage{iftex}
\RequirePDFTeX
\usepackage{cmap}
\usepackage[T2A]{fontenc}
\usepackage[utf8]{inputenc}
\usepackage[english,russian]{babel}
\usepackage{PTSerif}
\usepackage{amsmath,amssymb,microtype}
\usepackage[a4paper,left=23mm,right=23mm,top=20mm,bottom=21mm]{geometry}
\usepackage{xcolor,enumitem,fancyhdr,hyperref}

\definecolor{accent}{HTML}{187AA5}
\definecolor{ink}{HTML}{302B28}
\definecolor{muted}{HTML}{6D6560}
\definecolor{tint}{HTML}{EDF4F6}
\definecolor{rulecolor}{HTML}{A9CADA}
\hypersetup{unicode=true,colorlinks=true,urlcolor=accent,linkcolor=accent,
  pdftitle={Домашнее задание 3. Околокруговое приближение},
  pdfauthor={Чиняев Владимир Николаевич}}
\urlstyle{same}
\color{ink}
\setlength{\parindent}{0pt}
\setlength{\parskip}{3pt}
\linespread{1.08}
\setlist[enumerate]{leftmargin=*,label=\textcolor{accent}{\arabic*.},
  itemsep=3pt,topsep=4pt,parsep=0pt}
\setlist[itemize]{leftmargin=*,itemsep=3pt,topsep=3pt,parsep=0pt}
\setlength{\emergencystretch}{1em}
\pagestyle{fancy}
\fancyhf{}
\renewcommand{\headrulewidth}{0pt}
\renewcommand{\footrulewidth}{0.4pt}
\renewcommand{\footrule}{\hbox to\headwidth{\color{rulecolor}\leaders\hrule height\footrulewidth\hfill}}
\fancyfoot[L]{\footnotesize\color{muted}Введение в маневрирование КА\quad /\quad ДЗ 3}
\fancyfoot[R]{\footnotesize\color{muted}\thepage}

\newcommand{\taskheading}[2]{%
  \par\vspace{5pt}{\large\bfseries\color{accent}Задание #1. #2}\par\smallskip}


\newcommand{\hint}[1]{%
  \par\nobreak\smallskip\noindent
  \begingroup
  \setlength{\fboxsep}{7pt}%
  \setlength{\fboxrule}{0.4pt}%
  \fcolorbox{rulecolor}{tint}{%
    \begin{minipage}{\dimexpr\linewidth-2\fboxsep-2\fboxrule\relax}
    \small\textbf{\color{accent}Подсказка.} #1\par
    \end{minipage}}%
  \endgroup\par\smallskip}

\begin{document}

{\small\color{accent}ВВЕДЕНИЕ В МАНЕВРИРОВАНИЕ КА}\par
\vspace{3pt}
{\huge Домашнее задание 3}\par
{\large Околокруговое приближение}\par
{\small\color{muted}К лекции 4}
\par{\small\bfseries\color{accent}Срок сдачи: 19 октября 2026 года, 23:59}
\vspace{4pt}
{\color{rulecolor}\hrule}
\vspace{5pt}

Исследуйте пассивное движение КА в окрестности круговой орбиты.
Используйте уравнения первого приближения, полученные на лекции.
Оформите решение всех трёх заданий в предоставленном Jupyter Notebook.

\taskheading{1}{Чтение методички}

В методичке \texttt{near-circular-transfer.pdf} прочитайте
разделы 1.1--1.4 и 1.8 главы 1 «Уравнения движения»:
выбор координат, линеаризация, решение на пассивном участке,
параметры орбиты и границы применимости модели.

Дайте краткие обоснованные ответы:
\begin{enumerate}
\item Чем отличаются отклонение трансверсальной скорости $\Delta V_t$
и величина $r_0\Delta\dot u$? Запишите связь между ними в первом приближении.
\hint{Начните с кинематического равенства $V_t=r\dot u$ в разделе 1.1.
Подставьте $r=r_0+\Delta r$, $\dot u=n_0+\Delta\dot u$ и $V_t=V_0+\Delta V_t$;
раскройте скобки и отбросьте произведение малых отклонений.
Используйте $V_0=n_0r_0$ и поясните физический смысл полученных слагаемых.}
\item Какие величины должны оставаться малыми для применимости
околокругового приближения? Почему малость начальных отклонений
сама по себе не гарантирует применимости модели на всём интервале движения?
\hint{Выпишите безразмерные величины из (1.6), раздел 1.2.
Затем рассмотрите (1.12), раздел 1.3: найдите член, растущий со временем
при $a\ne r_0$, и сопоставьте его с условием малости $\Delta u$.
Объяснение границ применимости приведено в разделе 1.8.}
\end{enumerate}

\begingroup
\setlength{\fboxsep}{9pt}
\colorbox{tint}{\begin{minipage}{\dimexpr\linewidth-2\fboxsep\relax}
\textbf{Общие исходные данные}
\par\smallskip
Опорный аппарат движется по круговой орбите радиуса $r_0$.
Второй аппарат движется в той же плоскости и в том же направлении.
Управление и возмущения отсутствуют: $S=T=W=0$, $z=V_z=0$.
\[
 \mu=398600{,}44158\ \text{км}^3/\text{с}^2,\qquad r_0=7000\ \text{км}.
\]
\[
 V_0=\sqrt{\frac{\mu}{r_0}},\qquad n_0=\frac{V_0}{r_0},\qquad
 T_0=\frac{2\pi}{n_0},\qquad \varphi(t)=n_0t.
\]
В начальный момент аппараты находятся на одном радиальном луче:
\[
 \Delta r(0)=100\ \text{м},\qquad
 \Delta V_r(0)=0,\qquad \Delta u(0)=0.
\]
$\Delta r=r-r_0$, $\Delta V_t=V_t-V_0$, $\Delta u=u-\varphi$.
Координата $\eta=r_0\Delta u$ описывает продольное смещение
в первом приближении. Здесь $\Delta V_r,\Delta V_t$ --- отклонения
компонент скорости от кругового движения.\par
\end{minipage}}
\endgroup

В расчётах используйте километры, секунды и радианы.
В таблицах и на рисунках отклонения координат приводите в метрах,
отклонения скоростей --- в м/с.

\newpage
{\small\color{muted}Домашнее задание 3\hfill К лекции 4}\par
\vspace{5pt}

\taskheading{2}{Равенство начальных угловых скоростей}

В начальный момент угловые скорости аппаратов равны:
\[
 \dot u(0)=n_0,\qquad \Delta\dot u(0)=0.
\]

\begin{enumerate}
\item Найдите $\Delta V_t(0)$ и определите $a,e_x,e_y$ второго аппарата
в первом приближении. Сначала получите аналитические выражения,
затем вычислите численные значения.
\hint{Используйте связь скоростей из пункта 1 задания 1 и условие $\Delta\dot u(0)=0$.
Для $a,e_x,e_y$ возьмите (1.13) и (1.15) из раздела 1.4;
при $t=0$ имеем $\varphi=0$, поэтому $\cos\varphi=1$, $\sin\varphi=0$.
Нормировки даны в (1.9). Числа подставляйте после упрощения формул,
переведя $100$ м в $0{,}1$ км.}
\item Подставьте найденные параметры в решение пассивного движения
из методички. Получите выражения для $\Delta r(t)$ и $\Delta u(t)$.
Постройте траекторию на интервале $0\le t\le2T_0$ в координатах
$(\eta,\Delta r)$. По горизонтальной оси отложите $\eta$,
по вертикальной --- $\Delta r$; положительное направление $\eta$
совпадает с направлением движения опорного аппарата.
Отметьте положения при $t=0,T_0,2T_0$.
\hint{Возьмите выражение для $\Delta\widehat r$ из (1.11) и для $\Delta u$ из (1.12),
раздел 1.3. Подставьте найденные $a,e_x,e_y$ и положите
$\varphi_*=0$, $\Delta u(\varphi_*)=0$, $\varphi=n_0t$.
Вернитесь к размерной координате $\Delta r=r_0\Delta\widehat r$,
затем найдите $\eta=r_0\Delta u$. Проверьте формулы при $t=0$;
график строит готовая ячейка ноутбука.}
\item Составьте таблицу $\Delta r,\eta,\Delta V_r,\Delta V_t$
при $t=0,T_0,2T_0$.
Сохраняется ли равенство угловых скоростей при дальнейшем движении?
Объясните знак и накопление продольного смещения,
связав их с большой полуосью орбиты второго аппарата.
\hint{Для таблицы подставьте $\varphi=0,2\pi,4\pi$ в (1.11), (1.12);
скорости переведите из нормированного вида по (1.9).
Найдите $\Delta\dot u(t)$ по (1.7) или дифференцированием своего ответа.
По формуле $\Delta u(t)$ определите знак накопленного смещения;
его физический смысл объясните через среднее движение $n=\sqrt{\mu/a^3}$.}
\end{enumerate}

\vspace{4pt}
\textbf{Работа с ноутбуком}

В файле \texttt{hw\_3.ipynb} заданы исходные данные и подготовлены
функции расчёта пассивного движения по $a,e_x,e_y$,
построения траекторий и вывода таблиц.
Для каждого пункта оставлены текстовые ячейки для вывода формул
и пояснений, а также кодовые ячейки для численных результатов.
Заполните параметры орбит в заданиях 2 и 3 найденными вами значениями.
Готовые функции используют решение из методички;
вывод формул для заданных начальных условий выполните самостоятельно.

\newpage
{\small\color{muted}Домашнее задание 3\hfill К лекции 4}\par
\vspace{5pt}

\taskheading{3}{Периодическое движение}

Сохраните заданные $\Delta r(0),\Delta V_r(0),\Delta u(0)$,
а $\Delta V_t(0)$ выберите из условия периодичности движения
относительно опорного аппарата.

\begin{enumerate}
\item Найдите значение $\Delta V_t(0)$, при котором в линейной модели
состояние $\Delta r,\Delta u,\Delta V_r,\Delta V_t$
повторяется через период $T_0$.
Обоснуйте условие периодичности по выражению для $\Delta u(t)$
и определите соответствующие $a,e_x,e_y$.
\hint{В формуле (1.12) сравните $\Delta u(\varphi+2\pi)$ и $\Delta u(\varphi)$.
Тригонометрические члены повторяются; требование нулевого приращения
за период даст условие на $a$. Подставьте его в (1.13) и найдите
$\Delta V_t(0)$. Затем вычислите $e_x,e_y$ по (1.15) при $\varphi=0$.}
\item Получите выражения для $\Delta r(t)$ и $\Delta u(t)$.
Постройте траектории из заданий 2 и 3 на одном рисунке в двух панелях,
с одинаковым масштабом по осям $\eta$ и $\Delta r$ внутри каждой панели.
Отметьте положения при $t=0,T_0,2T_0$.
Дополните таблицу результатами задания 3 и проверьте
повторение всех четырёх компонент состояния через $T_0$.
\hint{Повторите подстановку в (1.11), (1.12) с параметрами задания 3.
Проверьте каждую компоненту состояния при $\varphi=0$ и $2\pi$,
в том числе $\Delta u$. Используйте готовые ячейки с двумя графиками,
таблицами и проверкой периодичности. Малые ненулевые разности
сопоставьте с погрешностью округления.}
\item Сравните результаты заданий 2 и 3. Объясните, почему равенство начальных
угловых скоростей и периодичность движения соответствуют
разным значениям $\Delta V_t(0)$.
Какую роль в этом различии играет большая полуось?
\hint{Сопоставьте найденные $\Delta V_t(0)$ и $a$ в двух заданиях.
Объясните, какое условие относится только к начальному моменту,
а какое --- к движению за целый период. Свяжите различие
с наличием или отсутствием линейного по $\varphi$ члена в (1.12).}
\end{enumerate}

\vspace{5pt}
{\color{rulecolor}\hrule}
\vspace{5pt}
\textbf{Что представить}

Сдайте один ноутбук с ответами на вопросы задания 1
и полными решениями заданий 2 и 3: аналитическими выражениями,
численными значениями начальных скоростей и параметров орбит,
траекториями, таблицами и объяснением результатов.
Формулы и рассуждения оформите в текстовых ячейках,
расчёты и рисунки --- в кодовых.
Перед отправкой перезапустите ядро, выполните все ячейки по порядку
и сохраните ноутбук с результатами вычислений и графиками.

Название файла: \texttt{hw\_3\_name\_surname.ipynb}.
Вместо \texttt{name\_surname} укажите свои имя и фамилию латиницей.

\end{document}
